\section{Identity Control Plane: State, Policy y Lifecycle}

La sección anterior explicó el wedge; esta entra en el sistema propiamente dicho. El identity control plane reúne HR, directory, customer signup, partner federation y service identity, que están dispersos, en un principal model unificado, y después aplica a cada acceso las políticas de authentication, authorization, session y lifecycle. No es la única truth de todos los datos, pero debe saber qué source tiene autoridad sobre qué attribute.

\subsection{El directory es la vista unificada del estado de identidad}

El \term{directory} (directorio) guarda el principal profile, los group, las relationship, el credential enrollment y el lifecycle state. La durable idea de Okta Universal Directory consiste en mapear las identity heterogéneas de Active Directory, LDAP, HR, las aplicaciones y las API a un modelo de datos predecible. La \term{source of truth} (fuente de autoridad) debe definirse por campo: por ejemplo, HR decide el employment status, la aplicación decide el role local y el IdP decide el authenticator enrollment; no se puede usar «el último en escribir» para sobrescribir toda la semántica.

\lecturefigure{02-identity-control-plane.png}{El identity control plane conecta las fuentes de identidad, el directorio unificado, el motor de políticas, las aplicaciones y la auditoría.}{Subtítulos locales 03:00--06:10; concepto oficial de Okta Universal Directory; redibujo conceptual.}

\begin{knowledgebox}{Lectura de la figura: primero pregunta de dónde viene el estado y adónde va}
La identity source se encarga de crear o actualizar los hechos; el directory normaliza principal, group y attribute; el policy engine toma decisiones de acceso según la aplicación de destino y el contexto; las apps y el audit reciben token, provisioning event y registros del sistema. Lo verdaderamente difícil son el conflict, la freshness, la deletion y la ownership: quién prevalece cuando HR y el directorio local están en conflicto, cuánto tarda en propagarse un evento de baja, cómo se recalcula un derived group y cómo demuestra la auditoría qué versión del estado usó un allow determinado.
\end{knowledgebox}

\subsection{Authentication, authorization, provisioning y audit}

La \term{multi-factor authentication} (autenticación multifactor, MFA) exige evidencias de distintas factor category, por ejemplo possession e inherence; \term{WebAuthn/passkey} usa credenciales de clave pública y puede ofrecer phishing-resistant authentication. Que la autenticación tenga éxito no significa que el usuario pueda acceder a todos los recursos: la authorization todavía debe revisar group, role, attribute, device, risk y requested action.

\lecturefigure{03-identity-lifecycle.png}{La identity infrastructure abarca cuatro tipos de problemas: autenticación, autorización, aprovisionamiento/revocación, y session/audit.}{Conceptos de Okta/OIDC/SCIM/WebAuthn; redibujo conceptual.}

\begin{knowledgebox}{Lectura de la figura: Login y lifecycle operan en escalas de tiempo distintas}
La authentication responde en segundos quién es el sujeto actual; la authorization restringe los permisos en cada acción sobre un recurso; el \term{provisioning} (aprovisionamiento) crea cuentas y entitlement en las aplicaciones posteriores; \term{SCIM} (System for Cross-domain Identity Management) sincroniza usuarios y group mediante una API estándar; session/audit, en cambio, se extiende hasta el cierre de sesión, un cambio de riesgo, la baja o una investigación. El SSO puede funcionar con normalidad y aun así una falla de offboarding deja una orphan account; por eso, los indicadores de seguridad deben cubrir todo el lifecycle.
\end{knowledgebox}

\begin{warningbox}{La MFA no es la solución universal contra el session theft}
La MFA reduce de forma notable el riesgo de que alguien se apodere de una cuenta solo con la contraseña, pero si el atacante roba una bearer session ya establecida, induce al usuario a aprobar un challenge equivocado o controla el proceso de recuperación o de help-desk, todavía puede eludir las garantías esperadas. El sistema necesita phishing-resistant authenticator, session binding, risk re-evaluation, recovery control y revoke rápido, en lugar de limitarse a contar el MFA enrollment.
\end{warningbox}

\subsection{Resumen de la sección}

Lo que unifica el identity control plane es la interfaz de estado y de decisión, no la eliminación de toda la source ownership. La autenticación, la autorización, el provisioning, la session y el audit pertenecen a ciclos de vida distintos, y para cada uno hay que definir por separado la correctness y la recovery.
